\section*{Hoofdstuk \ref{ch:b3:stem-cells} --- Stamcellen, regeneratie en veroudering}

\begin{solution}{exo:b3:stem-cells:1}
\omterm{def:b3:stem-cells:stem}{Zelfvernieuwing}: een deling die ten minste één dochter oplevert die nog
\omterm{def:b3:stem-cells:stem}{stamcel} is --- de \emph{Lgr5}-cellen aan de basis van de crypte.
Potentie: het bereik van lotgevallen dat een cel nog kan aannemen --- de
\omterm{def:b3:stem-cells:stem}{stamcel} van de crypte is multipotent en maakt opnemende cellen,
slijmbekercellen, entero-endocriene cellen, \omterm{prop:b3:stem-cells:crypt}{Panethcellen} en tuftcellen.
\omterm{def:b3:stem-cells:stem}{Nis}: de \omterm{prop:b3:stem-cells:crypt}{Panethcellen} en hun signalen van Wnt, Notch en EGF, die de
stamcelaard op haar plaats houden. \omterm{def:b3:stem-cells:stem}{Doorgangsversterkende cel}: een
vastgelegde dochter die in de crypte nog vier of vijf maal deelt voordat
zij op weg naar boven langs de villus differentieert.
\end{solution}

\begin{solution}{exo:b3:stem-cells:2}
Afzonderlijke mergcellen brachten kolonies op de milt voort met rode
cellen, granulocyten en bloedplaatjes: één cel, verscheidene lijnen ---
multipotentie, en het lineaire verband met de dosis liet zien dat elke
kolonie van één cel kwam. Voor de \omterm{def:b3:stem-cells:stem}{zelfvernieuwing} was een tweede muis
nodig: cellen uit een kolonie, in een tweede bestraalde ontvanger
gespoten, maakten nieuwe kolonies, dus had de stichtende cel dochters
met haar eigen vermogen voortgebracht.
\end{solution}

\begin{solution}{exo:b3:stem-cells:3}
\emph{Oct4}, \emph{Sox2}, \emph{Klf4}, \emph{c-Myc}. Het \omterm{def:b3:chromatin-epigenetics:nucleosome}{chromatine} van
de fibroblast heeft de genen van de \omterm{def:b3:stem-cells:pluripotency}{pluripotentie} achter \omterm{def:b3:chromatin-epigenetics:nucleosome}{heterochromatine}
en \omterm{def:b3:chromatin-epigenetics:methylation}{DNA-methylering} gesloten, zodat de factoren aanvankelijk weinig
toegankelijke plaatsen vinden en over vele celcycli
verbouwingsenzymen moeten aantrekken; de cel moet bovendien haar eigen
programma afzetten, aan de veroudering ontsnappen en een toevallige
tussentoestand doorlopen waarin de meeste cellen blijven steken of
sterven. Alleen de zeldzame cel waarin elke stap slaagt komt er weken
later pluripotent uit.
\end{solution}

\begin{solution}{exo:b3:stem-cells:4}
De planarie houdt pluripotente \omterm{def:b3:stem-cells:stem}{stamcellen} van de volwassene aan, de
\omterm{def:b3:stem-cells:regeneration}{neoblasten}, die naar de wond trekken en bouwen wat er ontbreekt. De
salamander heeft zo'n voorraad niet: gedifferentieerde cellen bij de
stomp dedifferentiëren tot \omterm{def:b3:stem-cells:stem}{voorlopercellen}, die met de \omterm{def:b3:stem-cells:stem}{stamcellen} van de
weefsels zelf een \omterm{def:b3:stem-cells:regeneration}{blasteem} vormen dat het ledemaatprogramma opnieuw
uitvoert. De transplantaties van Kragl met telkens één gemerkt weefsel
lieten zien dat de cellen van het \omterm{def:b3:stem-cells:regeneration}{blasteem} hun herkomst trouw blijven
--- spier maakt spier, kraakbeen kraakbeen --- zodat het \omterm{def:b3:stem-cells:regeneration}{blasteem} een
mengsel van tot hun lijn beperkte \omterm{def:b3:stem-cells:stem}{voorlopercellen} is en geen
pluripotente massa.
\end{solution}

\begin{solution}{exo:b3:stem-cells:5}
$(11 - 5)/0.075 = 80$ verdubbelingen. Vanaf \qty{8}{kb}:
$(8 - 5)/0.075 = 40$. Met een compensatie van \qty{60}{\%} is het netto
verlies \qty{30}{bp}: $6000/30
= 200$.
\end{solution}

\begin{solution}{exo:b3:stem-cells:6}
Veertien delingen van \omterm{def:b3:stem-cells:stem}{stamcellen} per dag leveren 28 dochters; er moeten
er 14 \omterm{def:b3:stem-cells:stem}{stamcel} blijven, dus leggen er 14 zich vast --- één vastgelegde
dochter per \omterm{def:b3:stem-cells:stem}{stamcel} per dag. Elk daarvan levert $2^{4} = 16$ cellen op:
$14\times 16 = 224$ villuscellen per dag, van de orde van de \num{300}
die worden waargenomen (een vijfde doorgangsdeling geeft \num{448}). De
laag van de \omterm{def:b3:stem-cells:stem}{stamcellen} draagt veertien cellen aan delingen bij, de
doorgangslaag de andere \num{200}.
\end{solution}

\begin{solution}{exo:b3:stem-cells:7}
$2\times 10^{11}/2^{20} = 1.9\times 10^{5}$ vastgelegde dochters per
dag; met \num{10000} \omterm{def:b3:stem-cells:stem}{stamcellen} zijn dat er 19 per \omterm{def:b3:stem-cells:stem}{stamcel} per dag ---
tegenover één deling per jaar. De telling van twintig delingen geldt per
lijn van \omterm{def:b3:stem-cells:stem}{stamcel} tot rode cel, maar de tussenliggende \omterm{def:b3:stem-cells:stem}{voorlopercellen}
zijn niet eenmalig: de voorraden multipotente en vastgelegde
\omterm{def:b3:stem-cells:stem}{voorlopercellen} houden zichzelf wekenlang in stand met hun eigen
delingen, zodat vrijwel alle deling in de lagen van de \omterm{def:b3:stem-cells:stem}{voorlopercellen}
gebeurt en er zelden een beroep op de \omterm{def:b3:stem-cells:stem}{stamcel} wordt gedaan.
\end{solution}

\begin{solution}{exo:b3:stem-cells:8}
$\mu(50) = 10^{-3}\mathrm{e}^{1.74} = 5.7\times 10^{-3}$; $\mu(70) =
10^{-3}\mathrm{e}^{3.48} = 0.032$; $\mu(90) = 10^{-3}\mathrm{e}^{5.22} =
0.19$ per jaar. Met $\mu_{0}/\gamma = 0.0115$: $S(70) = \exp[-0.0115
\times 31.5] = 0.70$, $S(90) = \exp[-0.0115\times 184] = 0.12$.
Mediaan: $30 + \ln(61)/0.087 = 77$. Halvering van $\mu_{0}$:
$\ln(121)/0.087 = 55$, mediaan 85 --- één verdubbelingstijd, acht jaar,
gewonnen. Halvering van $\gamma$ tot $0.0435$: $\ln(31)/0.0435 = 79$,
mediaan 109: de exponent vertragen wint vier maal zoveel als de
constante halveren.
\end{solution}

\begin{solution}{exo:b3:stem-cells:9}
Het overgebleven derde deel moet verdrievoudigen: $\log_{2}3 = 1.6$
delingen per cel. \omterm{def:b3:stem-cells:regeneration}{Compenserende groei}, omdat de bestaande kwabben
groter worden doordat de cellen ter plaatse delen; de weggenomen
kwabben, met hun gangen, vaten en bouw, worden niet herbouwd --- de
lever herwint haar massa en haar functie maar niet haar vorm.
\end{solution}

\begin{solution}{exo:b3:stem-cells:10}
Bij elke gebeurtenis gaat de omvang $k$ van de kloon met gelijke kans
naar $k + 1$ of $k - 1$, dus verandert haar verwachte omvang nooit;
eindigt de wandeling, dan is die $N$ (overname) of $0$ (verlies), en de
verwachting $N\,P + 0 = k$ geeft $P = k/N$. Voor één cel: $1/14$. Een
neutrale mutatie in één \omterm{def:b3:stem-cells:stem}{stamcel} wordt met kans $1/14$ in haar crypte
vastgelegd en gaat anders binnen maanden verloren --- de drift ruimt
dertien van de veertien mutaties op.
\end{solution}

\begin{solution}{exo:b3:stem-cells:11}
De wandeling is nu scheef: de kloon wint vaker een plaats dan zij er
een verliest, dus groeit haar verwachte omvang en nadert haar kans op
overname de zekerheid naarmate $k$ groeit --- bij een voordeel van
\qty{10}{\%} in een \omterm{def:b3:stem-cells:stem}{nis} van veertien ruwweg driemaal boven $1/14$ voor
één cel, en nog hoger zodra de kloon een paar plaatsen bezet. In de
loop van decennia nemen \omterm{def:b3:stem-cells:stem}{stamcellen} met zulke mutaties (\emph{DNMT3A},
\emph{TET2} in het merg) hun nissen over, en op zeventigjarige leeftijd
heeft een vijfde van de mensen een groot deel van het bloed uit één
kloon. Het is geen \omterm{def:b3:cancer-biology:hallmarks}{kanker}: de cellen differentiëren nog en
gehoorzamen hun hiërarchie. Maar de volgende mutatie heeft nu een kloon
van miljarden om in te ontstaan, en het risico op leukemie is
tienvoudig verhoogd.
\end{solution}

\begin{solution}{exo:b3:stem-cells:12}
Mediaan $= 30 + \gamma^{-1}\ln(1 + \gamma\ln 2/\mu_{0})$: voor $\mu_{0} =
10^{-2}$ is dat $30 + \ln 7/0.087 = 52$; voor $10^{-3}$ 77; voor
$10^{-4}$ $30 +
\ln 604/0.087 = 104$. In de zuivere vorm van Gompertz voegt elke
tienvoudige verlaging dezelfde $\ln 10/\gamma = 26$ jaar toe. In de
praktijk neemt de opbrengst af omdat de historische winst kwam van het
wegnemen van een leeftijdsonafhankelijke term --- besmetting,
bevalling, ongeval, de constante $A$ van Makeham in $\mu = A
+ \mu_{0}\mathrm{e}^{\gamma t}$ --- en zodra $A$ klein is ten opzichte
van de e-macht op de leeftijden die ertoe doen, verandert hem verder
verlagen bijna niets; de intrinsieke $\mu_{0}$ is veel minder
opgeschoven. Een verlaging van $\gamma$ met \qty{20}{\%} (tot $0.070$)
geeft $30 + \ln 49/0.070 = 86$: negen jaar, gewonnen op elke leeftijd
in plaats van door een paar doodsoorzaken uit te stellen.
\end{solution}

\begin{solution}{pb:b3:stem-cells:1}
\textbf{1.} $3.5\times 10^{11}$ per dag, $4\times 10^{6}$ per seconde.
\textbf{2.} $2^{20} \approx 10^{6}$ cellen per vastgelegde dochter;
$3.5\times 10^{11}/10^{6} = 3.3\times 10^{5}$ vastgelegde dochters per dag.
\textbf{3.} $33$ per \omterm{def:b3:stem-cells:stem}{stamcel} per dag, \num{12000} per jaar.
\textbf{4.} De voorraden \omterm{def:b3:stem-cells:stem}{voorlopercellen} moeten zichzelf in stand
houden: multipotente en vastgelegde \omterm{def:b3:stem-cells:stem}{voorlopercellen} delen wekenlang en
houden hun eigen aantallen op peil, zodat bijna elke deling van de
twintig plaatsvindt in lagen die zichzelf vernieuwen, en de \omterm{def:b3:stem-cells:stem}{stamcel}
zelden een vastgelegde dochter bijdraagt --- in de orde van eens per
jaar.
\textbf{5.} $10 - 20\times 0.06 = \qty{8.8}{kb}$: ruim boven
$L_{\text{crit}}$, en de rode cel deelt nooit meer; de reserve telt
alleen voor \omterm{def:b3:stem-cells:stem}{voorlopercellen} die opnieuw worden gebruikt.
\textbf{6.} Netto verlies $0.3\times 60 = \qty{18}{bp}$; $6000/18 = 333$
delingen; bij één per jaar $333$ jaar --- ver voorbij een mensenleven,
dus begrenzen de telomeren de \omterm{def:b3:stem-cells:stem}{stamcel} zelf niet; andere schade doet dat.
\textbf{7.} $14$ delingen per dag; $7$ vervangen verloren \omterm{def:b3:stem-cells:stem}{stamcellen} en
$7$ leggen zich vast; $7\times 2^{4} = 112$ villuscellen per dag.
\textbf{8.} $N^{2} = 196$ gebeurtenissen per cel bij één per dag:
ongeveer \qty{200}{dagen}, zes tot zeven maanden --- de tijd waarin de
crypten met confetti eenkleurig werden.
\textbf{9.} $1/14 \approx \qty{7}{\%}$; $10^{7}/14 \approx 7\times 10^{5}$
crypten.
\textbf{10.} De drift werpt dertien van elke veertien mutaties binnen
maanden weg. Levenslange \omterm{def:b3:stem-cells:stem}{asymmetrische deling} zou elke gemuteerde
\omterm{def:b3:stem-cells:stem}{stamcel} een leven lang behouden, zodat mutaties zich in elke lijn
zouden ophopen zonder ooit te worden opgeruimd.
\textbf{11.} De stamcelaard is een plaats in de \omterm{def:b3:stem-cells:stem}{nis} en een stel
signalen, geen blijvende eigenschap van een cel: elke cel die de
signalen van de \omterm{prop:b3:stem-cells:crypt}{Panethcellen} bereikt kan de rol op zich nemen.
\textbf{12.} Wek op één ogenblik een meerkleurige merking in de
\omterm{def:b3:stem-cells:stem}{stamcellen} op en volg de crypten. Symmetrische drift: de crypten worden
binnen maanden eenkleurig. \omterm{def:b3:stem-cells:stem}{Asymmetrische deling}: elke crypte behoudt al
haar kleuren onbeperkt, elk als een dun lint langs de villus.
\textbf{13.} $6000/60 = 100$ verdubbelingen.
\textbf{14.} Ongeveer $75$ over: de cel telt delingen, geen kalendertijd,
want het decennium in de vriezer kostte niets.
\textbf{15.} $(8000 - 4000)/30 = 133$ jaar na het twintigste --- op
leeftijd 153. De gemiddelde telomeerlengte van leukocyten begrenst het
leven dus niet zelf; de kortste telomeren, en de andere kenmerken,
tellen zwaarder.
\textbf{16.} \omterm{def:b3:cancer-biology:telomerase}{Telomerase} neemt één hindernis weg, de \omterm{thm:b3:stem-cells:hayflick}{replicatieve
veroudering}, en wordt daarom in vrijwel elke \omterm{def:b3:cancer-biology:hallmarks}{kanker} gevonden, die
onbeperkte deling nodig heeft; maar de fibroblasten die \omterm{def:b3:cancer-biology:telomerase}{telomerase}
kregen behielden hun controlepunten, hun contactremming en een gaaf
p53, zodat onbeperkte deling alleen er geen \omterm{def:b3:cancer-biology:hallmarks}{kankers} van maakte.
Noodzakelijk, niet voldoende.
\textbf{17.} $\mu(40) = 2.4\times 10^{-3}$, $\mu(60) = 0.014$, $\mu(80) =
0.077$, $\mu(100) = 0.44$ per jaar. $S = \exp[-0.0115(\mathrm{e}^{\gamma
(t-30)} - 1)]$: $0.98$ op 40, $0.87$ op 60, $0.42$ op 80 en $0.006$ op
100.
\textbf{18.} Mediaan $30 + \ln 61/0.087 = 77$. Tien procent leeft nog
bij $\mathrm{e}^{\gamma(t-30)} = 1 + \gamma\ln 10/\mu_{0} = 201$: $t = 30 +
5.3/0.087 = 91$.
\textbf{19.} \qty{0.67}{mm/dag}; $\log_{2}1000 = 10$ verdubbelingen in
60 dagen, één per zes dagen.
\textbf{20.} Ongeveer $10^{8}$ delingen, dus $0.1$ oncogene mutatie per
\omterm{def:b3:stem-cells:regeneration}{regeneratie} --- één op de tien ledematen. Axolotls krijgen vrijwel nooit
\omterm{def:b3:cancer-biology:hallmarks}{tumoren}: zij zijn koudbloedig en traag, met een stevige onderdrukking
van \omterm{def:b3:cancer-biology:hallmarks}{tumoren}, en de cellen van het \omterm{def:b3:stem-cells:regeneration}{blasteem} blijven tot hun lijn beperkt
en differentiëren opnieuw in plaats van te blijven delen; regenererend
weefsel onderdrukt zelfs opgewekte \omterm{def:b3:cancer-biology:hallmarks}{tumoren}. Een zoogdier heeft meer
cellen, een hogere temperatuur en stofwisseling, en een langer leven
waarin een ontsnapte kloon kan uitgroeien.
\textbf{21.} Haal de zenuwen uit de stomp en breng een bolletje aan dat
het kandidaat-eiwit afgeeft, of plaats weefsel dat door zenuwen is
voorbehandeld: keert de \omterm{def:b3:stem-cells:regeneration}{regeneratie} terug, dan is er een diffundeerbare
stof. Kandidaten die bij watersalamanders en axolotls zijn gevonden:
het eiwit nAG van de voorste gradiënt, neureguline-1 en FGF's.
\textbf{22.} Het Wnt-signaal is hoog aan de staart en laag aan de kop,
en de wond aan het voorste eind van een stuk maakt gewoonlijk een kop
omdat Wnt daar laag is; een remmer van Wnt uitschakelen verhoogt Wnt
overal, lezen beide wonden ``achterkant'', en groeien er twee staarten.
Het stuk kent zijn eigen as aan de gradiënt die in zijn spier is
achtergebleven, en de wond bouwt de gradiënt daarop weer op.
\textbf{23.} Het overgebleven derde deel moet verdrievoudigen: elke
levercel deelt één maal (tot twee derde) en de helft deelt nog eens ---
gemiddeld $1.6$ delingen. De overgebleven kwabben zwellen tot de
oorspronkelijke massa; de ontbrekende kwabben en hun gangen en vaten
worden niet herbouwd.
\textbf{24.} \omterm{def:b3:stem-cells:regeneration}{Regeneratie} die versleten en beschadigd weefsel vervangt
zou de ophoping vertragen die de exponent meet, wat $\gamma$ verlaagt en
de kromme afvlakt --- axolotls vertonen weinig veroudering. De prijs zou
een hogere constante term zijn door \omterm{def:b3:cancer-biology:hallmarks}{kankers} in een warm, groot,
langlevend lichaam.
\textbf{25.} Het aantal van Hayflick is ongeveer $100$ verdubbelingen;
een \omterm{def:b3:stem-cells:stem}{stamcel} met hulp van \omterm{def:b3:cancer-biology:telomerase}{telomerase} heeft zo'n $330$ delingen, meer dan
drie eeuwen bij één per jaar; de mediane levensduur is $77$ jaar.
\end{solution}
